\section{前期のまとめと試験}
\label{b15_summary_exam}

\subsection{授業内容}
\subsubsection{概要}
これまでの14回の授業を通じて，心理統計の基礎となる概念と手法について学んできた。本コマでは前期の内容を総括するとともに，学習内容の理解度を確認するための試験を実施する。試験は授業時間内に行われ，これまでに学んだ記述統計，相関と回帰，確率分布，推測統計と検定などの内容から出題される。試験を通じて，心理統計の基本概念の理解と実践的な応用力が身についているかを確認する。

\subsubsection{コマ主題細目}
\begin{description}
	\item[前期の総括] 前期では「記述」と「因果と相関」を中心に学んできた。心理学における数量的アプローチの意義から始まり，データの可視化，記述統計量，相関と回帰分析，確率分布と統計的推論，そして帰無仮説検定まで，心理統計の基礎となる概念と手法を体系的に学習してきた。これらの知識は，後期の応用的な内容(分散分析，因子分析など)を学ぶための基盤となる。また，Rを用いた実習を通じて，統計解析の実践的スキルも身につけてきた。前期で学んだ内容を振り返り，後期の学習へとつなげる。
	      \begin{flushright}
		      $\to$ これまでの各回の内容を復習しておくこと。
	      \end{flushright}

	\item[学習到達度の確認試験] 前期で学んだ内容の理解度を確認するため，授業時間内に試験を実施する。試験はマークシート形式で行われ，A4用紙2枚までの自筆メモの持ち込みが許可される。テスト理論に基づいて等価性が担保された問題が出題され，前期の学習内容全般から幅広く出題される。特に，記述統計の基本概念，相関と回帰分析の解釈，確率分布の特性，推測統計と帰無仮説検定の考え方など，心理統計の基本的な理解が問われる。
	      \begin{flushright}
		      $\to$ 試験の詳細については，別途配布される「試験に関するお知らせ」を参照のこと。
	      \end{flushright}

	\item[後期の学習に向けて] 前期では主にデータの記述と基本的な統計的推論について学んだ。後期ではより複雑な統計モデルや分析手法(分散分析，多変量解析など)について学び，実際の心理学研究に応用する方法を学ぶ。前期で習得した基礎知識をもとに，より高度な統計的思考と分析技術を身につけていく。特に，サンプルから母集団への一般化や，複数の要因が絡む実験計画の分析方法など，心理学研究の実践に直結する内容が中心となる。
	      \begin{flushright}
		      $\to$ 夏休み中に前期の内容を復習し，後期の学習に備えること。
	      \end{flushright}

\end{description}
\subsubsection{キーワード}
\begin{itemize}
	\item 記述統計と推測統計
	\item 相関と因果
	\item 確率分布と統計的推論
	\item 帰無仮説検定
	\item 統計的思考の実践
\end{itemize}
\subsection{授業情報}
\paragraph{コマの展開方法}
前期内容の総括(約15分)の後，授業時間内試験(約60分)を実施する
\paragraph{標準シラバスにおける位置づけ}
科目番号5；心理学統計法；1.心理学で用いられる統計手法(総括)
\subsubsection{予習・復習課題}
\paragraph{予習}
前期の全授業内容を復習しておくこと。特に以下の点について理解を深めておくことが重要である：
\begin{enumerate}
\item 記述統計の基本(平均，分散，標準偏差など)とその意味
\item 相関係数と回帰分析の基本概念と解釈
\item 確率分布(二項分布，正規分布など)の特徴と用途
\item 標本分布と中心極限定理の考え方
\item 推測統計の基本概念(点推定，区間推定)
\item 帰無仮説検定の論理と手順
\end{enumerate}
また，授業で扱ったRの基本操作とデータ分析の方法についても振り返っておくこと。試験では自筆のメモ(A4用紙2枚まで，両面可)の持ち込みが許可されるので，重要な概念や公式をまとめておくと良い。

\paragraph{復習}
試験終了後は，出題された問題について再度考え，不明点があれば教科書や参考書，授業ノートで確認しておくこと。また，後期の学習に向けて，以下の点を意識しておくと良い：

\begin{enumerate}
\item 心理学研究における統計的思考の重要性
\item データに基づく客観的な推論の方法
\item 統計ソフトウェア(R)を用いた実践的な分析スキル
\item 心理学の諸問題を統計的に捉える視点
\end{enumerate}

夏休み期間中に前期の内容を総復習し，特に苦手な分野があれば重点的に学習しておくことを推奨する。後期ではより高度な統計手法を学ぶため，前期の基礎的な内容をしっかりと理解しておくことが重要である。